% Options for packages loaded elsewhere
\PassOptionsToPackage{unicode}{hyperref}
\PassOptionsToPackage{hyphens}{url}
\documentclass[
  ignorenonframetext,
]{beamer}
\newif\ifbibliography
\usepackage{pgfpages}
\setbeamertemplate{caption}[numbered]
\setbeamertemplate{caption label separator}{: }
\setbeamercolor{caption name}{fg=normal text.fg}
\beamertemplatenavigationsymbolsempty
% remove section numbering
\setbeamertemplate{part page}{
  \centering
  \begin{beamercolorbox}[sep=16pt,center]{part title}
    \usebeamerfont{part title}\insertpart\par
  \end{beamercolorbox}
}
\setbeamertemplate{section page}{
  \centering
  \begin{beamercolorbox}[sep=12pt,center]{section title}
    \usebeamerfont{section title}\insertsection\par
  \end{beamercolorbox}
}
\setbeamertemplate{subsection page}{
  \centering
  \begin{beamercolorbox}[sep=8pt,center]{subsection title}
    \usebeamerfont{subsection title}\insertsubsection\par
  \end{beamercolorbox}
}
% Prevent slide breaks in the middle of a paragraph
\widowpenalties 1 10000
\raggedbottom
\AtBeginPart{
  \frame{\partpage}
}
\AtBeginSection{
  \ifbibliography
  \else
    \frame{\sectionpage}
  \fi
}
\AtBeginSubsection{
  \frame{\subsectionpage}
}
\usepackage{iftex}
\ifPDFTeX
  \usepackage[T1]{fontenc}
  \usepackage[utf8]{inputenc}
  \usepackage{textcomp} % provide euro and other symbols
\else % if luatex or xetex
  \usepackage{unicode-math} % this also loads fontspec
  \defaultfontfeatures{Scale=MatchLowercase}
  \defaultfontfeatures[\rmfamily]{Ligatures=TeX,Scale=1}
\fi
\usepackage{lmodern}
\ifPDFTeX\else
  % xetex/luatex font selection
\fi
% Use upquote if available, for straight quotes in verbatim environments
\IfFileExists{upquote.sty}{\usepackage{upquote}}{}
\IfFileExists{microtype.sty}{% use microtype if available
  \usepackage[]{microtype}
  \UseMicrotypeSet[protrusion]{basicmath} % disable protrusion for tt fonts
}{}
\makeatletter
\@ifundefined{KOMAClassName}{% if non-KOMA class
  \IfFileExists{parskip.sty}{%
    \usepackage{parskip}
  }{% else
    \setlength{\parindent}{0pt}
    \setlength{\parskip}{6pt plus 2pt minus 1pt}}
}{% if KOMA class
  \KOMAoptions{parskip=half}}
\makeatother
\setlength{\emergencystretch}{3em} % prevent overfull lines
\providecommand{\tightlist}{%
  \setlength{\itemsep}{0pt}\setlength{\parskip}{0pt}}
% Slide layout defaults for warning-clean scientific decks.
\usepackage{etoolbox}
\IfFileExists{xurl.sty}{\usepackage{xurl}}{}
\IfFileExists{seqsplit.sty}{\usepackage{seqsplit}}{\newcommand{\seqsplit}[1]{#1}}
\protected\def\breakseq#1{\seqsplit{#1}}
\protected\def\breaktt#1{\begingroup\ttfamily\seqsplit{#1}\endgroup}
\setlength{\emergencystretch}{6em}
\tolerance=5000
\hbadness=10000
\hfuzz=1pt
\setlength{\tabcolsep}{2pt}
\AtBeginEnvironment{longtable}{\tiny\renewcommand{\arraystretch}{0.86}\setlength{\tabcolsep}{1pt}}
\AtBeginEnvironment{tabular}{\tiny\renewcommand{\arraystretch}{0.86}\setlength{\tabcolsep}{1pt}}

% Natbib and cross-reference fallbacks — slides don't load natbib
% or cleveref, but combined-PDF manuscript prose may emit these
% commands. The fallback renders citations as a bracketed key list
% and cross-references as detokenized labels so slides don't fail on
% undefined control sequences or raw underscores. \providecommand is
% a no-op if packages load later.
\providecommand{\citep}[1]{[#1]}
\providecommand{\citet}[1]{#1}
\providecommand{\citealp}[1]{#1}
\providecommand{\citeauthor}[1]{#1}
\providecommand{\citeyear}[1]{#1}
\providecommand{\cref}[1]{\texttt{\detokenize{#1}}}
\providecommand{\Cref}[1]{\texttt{\detokenize{#1}}}
\usepackage{bookmark}
\IfFileExists{xurl.sty}{\usepackage{xurl}}{} % add URL line breaks if available
\urlstyle{same}
% fallback for those not using the hyperref driver hyperxmp:
\makeatletter
\@ifundefined{xmpquote}{\newcommand{\xmpquote}[1]{#1}}{}
\makeatother
\hypersetup{
  hidelinks,
  pdfcreator={LaTeX via pandoc}}

\author{\texorpdfstring{}{}}
\date{}

\begin{document}

\section{Introduction}\label{sec:introduction}

\begin{frame}[fragile,allowframebreaks]{Introduction}
Editorial review is one of the longest-lived bottlenecks in scientific
writing --- and an obvious target for the kind of \emph{reproducible
computational workflow} advocated by \breakseq{[@peng2011reproducible]}. A
working draft accumulates structural drift (skipped heading levels,
sections that have grown to 2,000 words while their siblings sit at
200), readability drift (a once-clean introduction now reads at the
Gunning Fog level of a legal contract \breakseq{[@gunning1952technique]}), and
citation drift (\texttt{\breakseq{[@key]}} references that no longer exist in
\texttt{references.bib}, or bibliography entries that nobody actually
cites). The traditional remedy --- a senior co-author re-reading the
manuscript with \breakseq{[@strunk2000elements]} in hand --- is slow,
non-reproducible, and impossible to apply continuously.

\breaktt{template\_prose\_project} exists to demonstrate that the
editorial-review pass can be expressed as a \textbf{deterministic,
configurable, infrastructure-backed pipeline}. This project carries no
novel research contribution of its own; its purpose is to show how to
compose existing template infrastructure into a complete, reproducible
editorial workflow.

The architecture is simple. \breaktt{manuscript/config.yaml} defines
policy: target grade-level band, citation-density floor,
heading-structure rules, bibliography-consistency policy.
\breaktt{src/pipeline/\_\_init\_\_.py::run\_prose\_pipeline} reads the
manuscript directory, calls
\href{../../../../infrastructure/prose/SKILL.md}{\breaktt{infrastructure.prose.analyze\_manuscript}}
to produce a \breaktt{ManuscriptReport}, cross-checks the cited keys
against
\href{../../../../infrastructure/reference/citation/SKILL.md}{\breaktt{infrastructure.reference.citation.parse\_bibfile}}
for the \texttt{references.bib}, evaluates each configured check, and
writes a markdown review report alongside JSON artefacts and three
figures. None of this is project-specific: a different project can
re-use the same infrastructure with a different \texttt{config.yaml} and
a different manuscript directory.

The remainder of this paper documents the methodology
(\breakseq{[@sec:methodology]}), the run-time results on the bundled manuscript
(\breakseq{[@sec:results]}), and the architectural lessons drawn from wiring
prose analysis through the template pipeline (\breakseq{[@sec:conclusion]}).
\end{frame}

\end{document}
